\item \textbf{Evaluación de reservas globales de uranio como combustible de fisión.}
Se estima que las reservas mundiales de uranio extraíbles a costo viable son de $4.40 \times 10^6\text{ toneladas métricas}$. El $0.700\%$ del uranio natural corresponde al isótopo fisionable $^{235}\text{U}$.
\begin{enumerate}
    \item Calcule la masa total en gramos del $^{235}\text{U}$ en estas reservas.
    \item Encuentre el número de moles de $^{235}\text{U}$ asociados.
    \item Determine el número total de núcleos de $^{235}\text{U}$.
    \item Asumiendo que se liberan $200\text{ MeV}$ netos en cada evento de fisión y que toda la energía se captura con éxito, calcule la energía total extraíble en Joules.
    \item Si el consumo de energía mundial se mantiene constante en $1.5 \times 10^{13}\text{ J/s}$, ¿durante cuántos años podrían estas reservas satisfacer la demanda energética global?
    \item Comente sobre la viabilidad a largo plazo de la energía nuclear de fisión convencional.
\end{enumerate}